\scanpage{172}\mspage{27}%
allant jusqu'au bout\unsure{,} de cette résistance à fleur de peau, on arrive à généraliser et à approfondir la notion de ``type d'homotopie'', et à apporter encore une nouvelle eau au développement d'un bon yoga homotopique. Voici où je veux en venir.

Parlons d'abord de faisceaux \inserted{(d'ens, ou de modules etc)}, pour simplifier, au lieu de champs, \inserted{et \unsure{plaçons-nous|plaçons nous} sur le topos étale d'un schéma $X$.} Les faisceaux loc. constants - modulo une condition de finitude supplémentaire assez anodine - forment les plus simples des faisceaux \uline{constructibles}, à la définition desquels ils servent de modèle. Supposant $X$ cohérent ($=$ qu. qct et qu. sép), les faisceaux constructibles généraux sont ceux pour lesquels il existe une partition \inserted{finie} $X=\bigcup\limits_{i\in I}X_i$\raisebox{-1.7ex}[0pt][0pt]{\makebox[0pt][r]{\parbox{1.5mm}{\figure{\begin{tikzpicture}[x=1mm,y=1mm]\fill (0,0) circle[radius=.14];\end{tikzpicture}}}\hspace{.8em}}} de $X$ en ``cellules'' \inserted{ou ``strates''} $X_i$, loc. fermées et constructibles, telles que \deleted{\unsure{chaque|chaq.}} la restriction de $F$ à chaque $X_i$ soit loc constant \deleted{et} (plus une condition de finitude\ldots) Ainsi la catégorie des faisceaux constructibles sur $X$ . (qui redonne la catégorie de tous les faisceaux en passant à une catégorie de ind-objets\ldots) \unsure{apparait|apparaît} elle-$\overline m$ comme limite inductive de catégories, associées à des partitions de plus en
% Small point at whole-image (1161,1361), below/right of the union index: reproduced by the zero-width raised box; no mathematical role is assigned.
% The separate point after X on L20 is reproduced as a physical point, not asserted to end a sentence.
% archive sheet number: 172
